%%%%%%%%%%%%%%%%%%%%%%%%%%%%%%%%%%%%%%%%%
% FRI Data Science_report LaTeX Template
% Version 1.0 (28/1/2020)
% 
% Jure Demšar (jure.demsar@fri.uni-lj.si)
%
% Based on MicromouseSymp article template by:
% Mathias Legrand (legrand.mathias@gmail.com) 
% With extensive modifications by:
% Antonio Valente (antonio.luis.valente@gmail.com)
%
% License:
% CC BY-NC-SA 3.0 (http://creativecommons.org/licenses/by-nc-sa/3.0/)
%
%%%%%%%%%%%%%%%%%%%%%%%%%%%%%%%%%%%%%%%%%


%----------------------------------------------------------------------------------------
%	PACKAGES AND OTHER DOCUMENT CONFIGURATIONS
%----------------------------------------------------------------------------------------
\documentclass[fleqn,moreauthors,10pt]{ds_report}
\usepackage[english]{babel}

\graphicspath{{fig/}}




%----------------------------------------------------------------------------------------
%	ARTICLE INFORMATION
%----------------------------------------------------------------------------------------

% Header
\JournalInfo{FRI Natural language processing course 2025/26}

% Interim or final report
\Archive{Project report} 
%\Archive{Final report} 

% Article title
\PaperTitle{Travel Planning Chatbot} 

% Authors (student competitors) and their info
\Authors{Nina Kokalj and Tia Šarkezi}

% Advisors
\affiliation{\textit{Advisors: Slavko Žitnik}}

% Keywords
\Keywords{Keyword1, Keyword2, Keyword3 ...}
\newcommand{\keywordname}{Keywords}


%----------------------------------------------------------------------------------------
%	ABSTRACT
%----------------------------------------------------------------------------------------

\Abstract{
The abstract goes here. The abstract goes here. The abstract goes here. The abstract goes here. The abstract goes here. The abstract goes here. The abstract goes here. The abstract goes here. The abstract goes here. The abstract goes here. The abstract goes here. The abstract goes here. The abstract goes here. The abstract goes here. The abstract goes here. The abstract goes here. The abstract goes here. The abstract goes here. The abstract goes here. The abstract goes here. The abstract goes here. The abstract goes here. The abstract goes here. The abstract goes here. The abstract goes here. The abstract goes here.
}

%----------------------------------------------------------------------------------------

\begin{document}

% Makes all text pages the same height
\flushbottom 

% Print the title and abstract box
\maketitle 

% Removes page numbering from the first page
\thispagestyle{empty} 

%----------------------------------------------------------------------------------------
%	ARTICLE CONTENTS
%----------------------------------------------------------------------------------------

\section*{Introduction}

Modern AI tourism has transitioned from static engines to dynamic, conversational agents. This project focuses on the development of a travel planning chatbot designed to take the stress out of vacation planning. By analyzing user preferences such as geographic region, seasonality, and climate, the bot provides tailored destination recommendations. The system generates curated itineraries and highlights most popular attractions, ensuring every suggestion aligns with the user's specific budget and travel style. To ensure the highest relevance, the chatbot utilizes a smart filtering and scoring system to rank destinations based on their ratings, distance and popularity with other users.

\subsection{Datasets}

To train the travel chatbot, we will use a combination of domain-specific knowledge datasets and dialogue-oriented datasets.

\paragraph{Travel Knowledge Datasets}
These datasets provide information about destinations, attractions, and general travel content:
\begin{itemize}
    \item \textbf{Wikivoyage} – a travel guide with detailed information about countries, cities, and points of interest, including things to do and travel tips. The data is available as an XML dump and needs to be parsed before use \cite{wikivoyage_dump}.
    
    \item \textbf{GeoNames} – a geographical database with information about locations such as cities, countries, and landmarks \cite{geonames}.
\end{itemize}

\paragraph{Dialogue Datasets}
These datasets help the chatbot learn how to communicate naturally in travel-related situations:
\begin{itemize}
    \item \textbf{MultiWOZ 2.2} – a widely used conversational dataset in JSON format, containing dialogues about booking hotels, restaurants, and transportation \cite{multiwoz_github}.
    
    \item \textbf{DSTC8} – a multi-domain dialogue dataset that includes travel-related conversations, such as hotel and restaurant scenarios \cite{dstc8_sgd_github}.
\end{itemize}

\paragraph{Review Datasets}
Review datasets include real user opinions and experiences, helping the chatbot give more realistic recommendations and understand sentiment:

\begin{itemize}
    \item \textbf{515K Hotel Reviews in Europe} – a large-scale dataset containing hotel reviews across Europe \cite{hotel_reviews_europe}.
    
    \item \textbf{TripAdvisor Hotel Reviews} – a dataset of 20,000 hotel reviews from TripAdvisor \cite{tripadvisor_reviews}.
    
    \item \textbf{Booking.com Reviews (30k)} – a dataset containing approximately 30,000 hotel reviews from Booking.com \cite{booking_reviews}.
\end{itemize}

\paragraph{Additional Datasets (Hugging Face)}
To further improve performance, a few travel-related datasets from Hugging Face will also be used:
\begin{itemize}
    \item \texttt{fathyshalab/clinic-travel} – contains travel-related conversations \cite{clinic_travel}.
    \item \texttt{osunlp/TravelPlanner} – designed for itinerary planning and travel assistance \cite{travelplanner}.
    \item \texttt{soniawmeyer/travel-conversations} – contains curated travel dialogues for fine-tuning \cite{travel_conversations}.
\end{itemize}

\subsection{Related work}

Studies \cite{Elangovan2025, travalaManyTravelers, karlovic2025context} demonstrate that while LLMs excel at dialogue, they struggle with logistical "hard constraints" like budgets and schedules. To solve this, researchers are now using hybrid architectures that pair LLMs with deterministic filtering and scoring systems. These hybrid models have achieved up to 91\% recommendation accuracy, significantly reducing decision fatigue for the 40\% of global travelers who now use AI as their primary trip-discovery tool. By integrating user-specific weights for travel style and budget, these systems deliver high-fidelity itineraries that mirror the expertise of a human travel consultant.

%------------------------------------------------

\section*{Methods}



\subsection*{Current model performance}

The \textbf{Llama-3.1-8B-Instruct} model demonstrated strong reasoning and instruction-following capabilities directly out-of-the-box, eliminating the need for expensive fine-tuning. However, its utility is limited by a knowledge cutoff in late 2023, making it unable to process real-time data or events from 2023 onwards. To solve this, we integrated the model into an agentic framework. By providing the model with API access, we effectively extended its capabilities beyond its training data, allowing it to retrieve live information and perform tasks involving current, dynamic data.


\subsection*{Agents}

To extend the capabilities of the chatbot beyond the knowledge stored in the language model, we implemented a set of \textbf{tool-using agents}. These agents are designed to query external data sources when necessary, enabling the system to provide more accurate and up-to-date information. All agents are based on a language model combined with the \texttt{CodeAgent} framework.


\subsubsection*{Train Agent}

The goal of this component was to extend the chatbot with the ability to access \textbf{up-to-date train schedules}, which are not reliably stored in the language model itself.

For retrieving train connections, we implemented a custom tool using the ÖBB (Austrian Federal Railways) HAFAS API. This tool resolves station names and queries journey data, returning structured information such as departure time, arrival time, and the full route with transfers.

While the integration proved effective, achieving reliable behavior required careful prompt engineering. In early experiments, the agent exhibited several issues: it repeatedly called the tool with different departure times, searched for connections in the past (e.g., starting at midnight for queries referring to ``today''), and sometimes returned incorrect results by selecting times from intermediate transfers instead of the actual departure from the origin.

To address these problems, we introduced explicit rules in the \textbf{system prompt}. The agent is instructed to call the transport tool exactly once, use the provided datetime (or the current time if none is specified), and base its answer on the first leg of the journey. Additionally, the final response must include the departure time from the origin, the arrival time at the destination, and the full route with transfers, instead of returning only a single time.


\subsubsection*{Weather Forecast Agent}
The system uses the \textbf{OpenWeatherMap} 5-Day/3-Hour Forecast API to retrieve granular data, including min/max temperatures, weather conditions, and precipitation probability. While the retrieval and parsing of the raw JSON data were successful and factually accurate, the model produced a semi-correct output due to a failure in constraint following. It reported the full five days of available tool data despite the user's explicit request for only a 3 day forecast, prioritizing a complete summary of the observation block over the prompt's temporal limitations.

\subsection*{Further Improvements}

Future development will focus on integrating APIs for \textbf{flights, accommodations, and local amenities} to provide real-time pricing, availability, and localized reviews for restaurants and attractions. Adding location-based services will allow the model to identify specific points of interest. Furthermore, a dedicated itinerary tool will synthesize weather, transit times, and opening hours into cohesive, hour-by-hour schedules, transitioning the agent from a basic information retriever into a proactive end-to-end trip planner.


\subsection*{Figures}

You can insert figures that span over the whole page, or over just a single column. The first one, \figurename~\ref{fig:column}, is an example of a figure that spans only across one of the two columns in the report.

\begin{figure}[ht]\centering
	\includegraphics[width=\linewidth]{single_column.pdf}
	\caption{\textbf{A random visualization.} This is an example of a figure that spans only across one of the two columns.}
	\label{fig:column}
\end{figure}

On the other hand, \figurename~\ref{fig:whole} is an example of a figure that spans across the whole page (across both columns) of the report.

% \begin{figure*} makes the figure take up the entire width of the page
\begin{figure*}[ht]\centering 
	\includegraphics[width=\linewidth]{whole_page.pdf}
	\caption{\textbf{Visualization of a Bayesian hierarchical model.} This is an example of a figure that spans the whole width of the report.}
	\label{fig:whole}
\end{figure*}


\subsection*{Tables}

Use the table environment to insert tables.

\begin{table}[hbt]
	\caption{Table of grades.}
	\centering
	\begin{tabular}{l l | r}
		\toprule
		\multicolumn{2}{c}{Name} \\
		\cmidrule(r){1-2}
		First name & Last Name & Grade \\
		\midrule
		John & Doe & $7.5$ \\
		Jane & Doe & $10$ \\
		Mike & Smith & $8$ \\
		\bottomrule
	\end{tabular}
	\label{tab:label}
\end{table}


\subsection*{Code examples}

You can also insert short code examples. You can specify them manually, or insert a whole file with code. Please avoid inserting long code snippets, advisors will have access to your repositories and can take a look at your code there. If necessary, you can use this technique to insert code (or pseudo code) of short algorithms that are crucial for the understanding of the manuscript.

\lstset{language=Python}
\lstset{caption={Insert code directly from a file.}}
\lstset{label={lst:code_file}}
\lstinputlisting[language=Python]{code/example.py}

\lstset{language=R}
\lstset{caption={Write the code you want to insert.}}
\lstset{label={lst:code_direct}}
\begin{lstlisting}
import(dplyr)
import(ggplot)

ggplot(diamonds,
	   aes(x=carat, y=price, color=cut)) +
  geom_point() +
  geom_smooth()
\end{lstlisting}

%------------------------------------------------

\section*{Results}

Use the results section to present the final results of your work. Present the results in a objective and scientific fashion. Use visualisations to convey your results in a clear and efficient manner. When comparing results between various techniques use appropriate statistical methodology.

\subsection*{More random text}

This text is inserted only to make this template look more like a proper report. Lorem ipsum dolor sit amet, consectetur adipiscing elit. Etiam blandit dictum facilisis. Lorem ipsum dolor sit amet, consectetur adipiscing elit. Interdum et malesuada fames ac ante ipsum primis in faucibus. Etiam convallis tellus velit, quis ornare ipsum aliquam id. Maecenas tempus mauris sit amet libero elementum eleifend. Nulla nunc orci, consectetur non consequat ac, consequat non nisl. Aenean vitae dui nec ex fringilla malesuada. Proin elit libero, faucibus eget neque quis, condimentum laoreet urna. Etiam at nunc quis felis pulvinar dignissim. Phasellus turpis turpis, vestibulum eget imperdiet in, molestie eget neque. Curabitur quis ante sed nunc varius dictum non quis nisl. Donec nec lobortis velit. Ut cursus, libero efficitur dictum imperdiet, odio mi fermentum dui, id vulputate metus velit sit amet risus. Nulla vel volutpat elit. Mauris ex erat, pulvinar ac accumsan sit amet, ultrices sit amet turpis.

Phasellus in ligula nunc. Vivamus sem lorem, malesuada sed pretium quis, varius convallis lectus. Quisque in risus nec lectus lobortis gravida non a sem. Quisque et vestibulum sem, vel mollis dolor. Nullam ante ex, scelerisque ac efficitur vel, rhoncus quis lectus. Pellentesque scelerisque efficitur purus in faucibus. Maecenas vestibulum vulputate nisl sed vestibulum. Nullam varius turpis in hendrerit posuere.

Nulla rhoncus tortor eget ipsum commodo lacinia sit amet eu urna. Cras maximus leo mauris, ac congue eros sollicitudin ac. Integer vel erat varius, scelerisque orci eu, tristique purus. Proin id leo quis ante pharetra suscipit et non magna. Morbi in volutpat erat. Vivamus sit amet libero eu lacus pulvinar pharetra sed at felis. Vivamus non nibh a orci viverra rhoncus sit amet ullamcorper sem. Ut nec tempor dui. Aliquam convallis vitae nisi ac volutpat. Nam accumsan, erat eget faucibus commodo, ligula dui cursus nisi, at laoreet odio augue id eros. Curabitur quis tellus eget nunc ornare auctor.


%------------------------------------------------

\section*{Discussion}

Use the Discussion section to objectively evaluate your work, do not just put praise on everything you did, be critical and exposes flaws and weaknesses of your solution. You can also explain what you would do differently if you would be able to start again and what upgrades could be done on the project in the future.


%------------------------------------------------

\section*{Acknowledgments}

Here you can thank other persons (advisors, colleagues ...) that contributed to the successful completion of your project.


%----------------------------------------------------------------------------------------
%	REFERENCE LIST
%----------------------------------------------------------------------------------------
\bibliographystyle{unsrt}
\bibliography{report}


\end{document}
